\documentclass{article}
\usepackage[utf8]{inputenc}
\usepackage{amsmath}
\usepackage{amssymb}
\usepackage{amsthm}
\usepackage{parskip}

\newtheorem{proposition}{Proposition}

\title{Multi-interval examples improving the one-interval bound}
\author{Ferran Espu\~na}
\date{\today}

\begin{document}

\maketitle

We write \(\mathbb T=\mathbb R/\mathbb Z\). A set \(A\subset \mathbb T\) is
\(k\)-sum-free if
\[
    0\notin A+A-kA.
\]
Equivalently, there are no \(x,y,z\in A\) with \(x+y=kz\) in \(\mathbb T\).

For comparison, suppose \(A=I\) is a single open interval of length \(\alpha\).
Then \(I+I-kI\) is an open interval of length \((k+2)\alpha\) on the real line.
If \((k+2)\alpha>1\), its image in \(\mathbb T\) is all of \(\mathbb T\). Thus
the trivial one-interval bound is
\[
    \alpha\leq \frac1{k+2}.
\]
The examples below exceed this bound for \(k=4\) and \(k=5\). We also record
a general equally spaced construction which gives examples of density larger
than \(1/8\) for every \(k\geq 3\).

We will use the following finite certificate. If
\[
    I_r=(a_r,b_r),
\]
then
\[
    I_i+I_j-kI_\ell
    =
    \bigl(a_i+a_j-kb_\ell,\; b_i+b_j-ka_\ell\bigr).
\]
It is therefore enough, for each \(i\leq j\) and each \(\ell\), to find an
integer \(n\) such that
\[
    n\leq a_i+a_j-kb_\ell
    \quad\text{and}\quad
    b_i+b_j-ka_\ell\leq n+1.
\]
Then the open interval \(I_i+I_j-kI_\ell\) contains no integer, and hence
contributes no point equal to \(0\) in \(\mathbb T\).

\begin{proposition}[An equally spaced construction]\label{prop:equal_spacing}
Let \(k\geq 3\), and let \(m\) be an integer satisfying
\[
    1\leq m < \frac{k+2}{2}.
\]
Set
\[
    \delta=\frac{k-2m+2}{k(k+2)}.
\]
Choose \(a\in\mathbb T\) such that
\[
    (2-k)a\equiv k\delta \pmod 1,
\]
and define
\[
    A_{k,m}
    =
    \bigcup_{r=0}^{m-1}
    \left(a+\frac r k,\; a+\frac r k+\delta\right)
    \subset \mathbb T.
\]
Then \(A_{k,m}\) is \(k\)-sum-free and
\[
    \mu(A_{k,m})
    =
    \frac{m(k-2m+2)}{k(k+2)}.
\]
In particular, optimizing over \(m\) gives the lower bound
\[
    \sup \{\mu(A): A\subset\mathbb T \text{ is } k\text{-sum-free}\}
    \geq
    \frac{\left\lfloor (k+2)^2/8\right\rfloor}{k(k+2)}
    >
    \frac18.
\]
\end{proposition}

\begin{proof}
The intervals are disjoint because \(\delta<1/k\). Write
\[
    I_r=\left(a+\frac r k,\; a+\frac r k+\delta\right),
    \qquad 0\leq r\leq m-1.
\]
For \(0\leq r,s,t\leq m-1\), the interval \(I_r+I_s-kI_t\) has endpoints
\begin{align*}
    L_{rst}
    &=
    (2-k)a+\frac{r+s}{k}-t-k\delta, \\
    R_{rst}
    &=
    (2-k)a+\frac{r+s}{k}-t+2\delta.
\end{align*}
By the choice of \(a\), there is an integer \(q\) such that
\[
    (2-k)a=q+k\delta.
\]
Thus
\[
    L_{rst}=q-t+\frac{r+s}{k}
\]
and
\[
    R_{rst}
    =
    q-t+\frac{r+s}{k}+(k+2)\delta
    =
    q-t+\frac{r+s}{k}+1-\frac{2m-2}{k}.
\]
Since \(0\leq r+s\leq 2m-2\), we have
\[
    q-t
    \leq L_{rst}
    <
    R_{rst}
    \leq q-t+1.
\]
Therefore each open interval \(I_r+I_s-kI_t\) contains no integer. Hence
\(0\notin A_{k,m}+A_{k,m}-kA_{k,m}\) in \(\mathbb T\).

The measure is \(m\delta\). It remains only to maximize
\[
    m(k-2m+2)=m((k+2)-2m)
\]
over integers \(m\). This quadratic has maximum
\[
    \left\lfloor \frac{(k+2)^2}{8}\right\rfloor.
\]
Writing \(u=k+2\), we also have
\[
    8\left\lfloor \frac{u^2}{8}\right\rfloor
    \geq u^2-7
    >
    u^2-2u
    =
    k(k+2),
\]
because \(u\geq 5\). This proves that the optimized density is always strictly
larger than \(1/8\).
\end{proof}

\begin{proposition}[A \(4\)-sum-free set of three open intervals]
Let
\[
    A_4=
    \left(\frac13,\frac5{12}\right)
    \cup
    \left(\frac{25}{48},\frac{101}{192}\right)
    \cup
    \left(\frac7{12},\frac23\right)
    \subset \mathbb T.
\]
Then \(A_4\) is \(4\)-sum-free and
\[
    \mu(A_4)
    =
    \frac1{12}+\frac1{192}+\frac1{12}
    =
    \frac{11}{64}
    >
    \frac16.
\]
Thus the one-interval bound for \(k=4\) is improved by \(1/192\).
\end{proposition}

\begin{proof}
Let the three intervals in \(A_4\) be \(I_1,I_2,I_3\), in the displayed order.
For each triple \(i\leq j\) and \(\ell\), the following table gives
\[
    L_{ij\ell}=a_i+a_j-4b_\ell,
    \qquad
    R_{ij\ell}=b_i+b_j-4a_\ell,
\]
and an integer \(n\) with \(n\leq L_{ij\ell}<R_{ij\ell}\leq n+1\).
\[
\begin{array}{c|c|c|c}
(i,j,\ell) & n & L_{ij\ell} & R_{ij\ell} \\
\hline
(1,1,1) & -1 & -1 & -\frac12 \\
(1,1,2) & -2 & -\frac{23}{16} & -\frac54 \\
(1,1,3) & -2 & -2 & -\frac32 \\
(1,2,1) & -1 & -\frac{13}{16} & -\frac{25}{64} \\
(1,2,2) & -2 & -\frac54 & -\frac{73}{64} \\
(1,2,3) & -2 & -\frac{29}{16} & -\frac{89}{64} \\
(1,3,1) & -1 & -\frac34 & -\frac14 \\
(1,3,2) & -2 & -\frac{19}{16} & -1 \\
(1,3,3) & -2 & -\frac74 & -\frac54 \\
(2,2,1) & -1 & -\frac58 & -\frac9{32} \\
(2,2,2) & -2 & -\frac{17}{16} & -\frac{33}{32} \\
(2,2,3) & -2 & -\frac{13}{8} & -\frac{41}{32} \\
(2,3,1) & -1 & -\frac9{16} & -\frac9{64} \\
(2,3,2) & -1 & -1 & -\frac{57}{64} \\
(2,3,3) & -2 & -\frac{25}{16} & -\frac{73}{64} \\
(3,3,1) & -1 & -\frac12 & 0 \\
(3,3,2) & -1 & -\frac{15}{16} & -\frac34 \\
(3,3,3) & -2 & -\frac32 & -1
\end{array}
\]
Since the intervals are open, equality with an integer endpoint in the table
does not introduce an integer point. Hence every interval
\(I_i+I_j-4I_\ell\) misses all integers, so \(0\notin A_4+A_4-4A_4\) in
\(\mathbb T\).
\end{proof}

\begin{proposition}[A \(5\)-sum-free set of two open intervals]
Let
\[
    A_5=
    \left(\frac2{35},\frac17\right)
    \cup
    \left(\frac67,\frac{33}{35}\right)
    \subset \mathbb T.
\]
Then \(A_5\) is \(5\)-sum-free and
\[
    \mu(A_5)
    =
    \frac3{35}+\frac3{35}
    =
    \frac6{35}
    >
    \frac17.
\]
Thus the one-interval bound for \(k=5\) is improved by \(1/35\).
\end{proposition}

\begin{proof}
The numerical experiment found a third interval of length \(0\), namely
\((6/7,6/7)\). We omit it here and keep only the two non-empty intervals above.
Let these intervals be \(I_1,I_2\).

For each triple \(i\leq j\) and \(\ell\), set
\[
    L_{ij\ell}=a_i+a_j-5b_\ell,
    \qquad
    R_{ij\ell}=b_i+b_j-5a_\ell.
\]
The following table gives an integer \(n\) with
\[
    n\leq L_{ij\ell}<R_{ij\ell}\leq n+1.
\]
\[
\begin{array}{c|c|c|c}
(i,j,\ell) & n & L_{ij\ell} & R_{ij\ell} \\
\hline
(1,1,1) & -1 & -\frac35 & 0 \\
(1,1,2) & -5 & -\frac{23}{5} & -4 \\
(1,2,1) & 0 & \frac15 & \frac45 \\
(1,2,2) & -4 & -\frac{19}{5} & -\frac{16}{5} \\
(2,2,1) & 1 & 1 & \frac85 \\
(2,2,2) & -3 & -3 & -\frac{12}{5}
\end{array}
\]
Again, the intervals are open, so the integer endpoints shown in the table are
not included. Therefore every interval \(I_i+I_j-5I_\ell\) misses all integers,
and \(0\notin A_5+A_5-5A_5\) in \(\mathbb T\).
\end{proof}

\end{document}
